% !TeX encoding = UTF-8
% !TeX TXS-program:compile = txs:///latexmk
%% =====================================================================
%%  aide-memoire.tex — Aide-mémoire d'UNE page : toutes les commandes.
%%  À imprimer ou garder ouvert pendant la rédaction.
%%  (\cmd~…~ affiche du code littéral ; le délimiteur ~ ne doit pas
%%   apparaître dans le code affiché.)
%% =====================================================================
\documentclass[a4paper,10pt]{scrartcl}
% !TeX encoding = UTF-8
%% =====================================================================
%%  style/preamble.tex — Packages, polices, mise en page, mathématiques,
%%  unités, figures/tableaux, code, bibliographie, hyperliens, titre.
%%
%%  NE PAS MODIFIER pour un devoir ordinaire :
%%    - identité permanente (nom, université, logo) → config.tex
%%    - informations du document (titre, cours, type) → début du fichier .tex
%%  Compilation : LuaLaTeX uniquement (latexmk s'en charge).
%%
%%  Ordre de chargement :
%%     \documentclass → style/preamble → style/macros → config → \docsetup{…}
%% =====================================================================

%% ---------------------------------------------------------------------
%%  0. Garde-fou : ce template exige LuaLaTeX
%% ---------------------------------------------------------------------
\RequirePackage{iftex}
\RequirePackage{etoolbox}
\ifluatex\else
  \PackageError{preamble}{Ce template doit être compilé avec LuaLaTeX}
    {Dans TeXstudio : Options > Configurer TeXstudio > Compilation > Compilateur par défaut = latexmk (voir README.md).}
\fi

%% ---------------------------------------------------------------------
%%  1. Palette : TEXTE ENTIÈREMENT NOIR ET GRIS (aucune couleur de thème)
%%     Les seules couleurs du document viennent des images importées, des
%%     courbes tracées et de la coloration syntaxique du code (plus bas).
%%     Pour réintroduire un accent coloré partout (liens, puces, filets,
%%     remarques…), il suffit de changer la définition de AccentColor.
%% ---------------------------------------------------------------------
\RequirePackage{xcolor}
\definecolor{TextColor}{gray}{0.08}      % texte : noir doux
\definecolor{MutedColor}{gray}{0.40}     % informations secondaires : gris foncé
\definecolor{RuleColor}{gray}{0.72}      % filets, grilles : gris clair
\definecolor{FillColor}{gray}{0.96}      % fonds très légers (code)
\colorlet{AccentColor}{TextColor}        % pas d'accent coloré : tout en noir
\colorlet{LinkColor}{TextColor}          % liens et citations : noir (cliquables)
\AtBeginDocument{\color{TextColor}}

%% Couleurs des DONNÉES et du CODE : indépendantes du thème ci-dessus.
%% Les figures importées (PDF, PNG…) ne sont jamais recolorées ni passées en gris.
\definecolor{DataBlue}{RGB}{31,119,180}  % courbes, points (convention matplotlib)
\definecolor{DataRed}{RGB}{214,39,40}
\definecolor{CodeKeyword}{RGB}{0,51,153}   % mots-clés : bleu
\definecolor{CodeBuiltin}{RGB}{0,115,115}  % fonctions prédéfinies : vert-bleu
\definecolor{CodeComment}{RGB}{85,125,85}  % commentaires : vert doux
\definecolor{CodeString}{RGB}{163,36,36}   % chaînes : rouge brique
\definecolor{CodeNumber}{RGB}{176,92,0}    % nombres : ocre
\definecolor{CodeLineNo}{gray}{0.55}

%% ---------------------------------------------------------------------
%%  2. Mathématiques : amsmath/mathtools AVANT unicode-math
%%     (amssymb et bm sont volontairement omis : unicode-math fournit tous
%%      les symboles AMS et les alphabets gras/ISO ; les charger ferait
%%      doublon, voire conflit.)
%% ---------------------------------------------------------------------
\RequirePackage{amsmath}
\RequirePackage{mathtools}
\RequirePackage{amsthm}

%% ---------------------------------------------------------------------
%%  3. Polices : Libertinus Serif (texte), Libertinus Math (mathématiques),
%%     Libertinus Sans (titres, étiquettes), Libertinus Mono (code), réglés pour
%%     LuaLaTeX selon test_fonts/CLAUDE-libertinus.md. fontspec + unicode-math.
%%     Chargées par NOM DE FICHIER : fonctionne sans que la base de polices de
%%     luaotfload soit à jour. Si les fichiers sont absents, repli sur Latin
%%     Modern avec un AVERTISSEMENT dans le log.
%% ---------------------------------------------------------------------
\RequirePackage{fontspec}
\RequirePackage[math-style=TeX,bold-style=TeX,warnings-off={mathtools-colon,mathtools-overbracket}]{unicode-math}

\IfFontExistsTF{LibertinusSerif-Regular.otf}
  {\IfFontExistsTF{LibertinusMath-Regular.otf}
    {\def\TemplateFonts{libertinus}}
    {\def\TemplateFonts{lm}}}
  {\def\TemplateFonts{lm}}

\ifdefstring{\TemplateFonts}{libertinus}{%
  % Réglage « Libertinus pour LuaLaTeX » (test_fonts/CLAUDE-libertinus.md) :
  % gras en Semibold (le Bold est trop lourd), maths avec intégrales inclinées (ss08).
  \setmainfont{LibertinusSerif}[
    Extension      = .otf,
    UprightFont    = *-Regular,
    ItalicFont     = *-Italic,
    BoldFont       = *-Semibold,
    BoldItalicFont = *-SemiboldItalic]
  \setsansfont{LibertinusSans}[
    Extension = .otf, UprightFont = *-Regular,
    ItalicFont = *-Italic, BoldFont = *-Bold]
  % Code (programmes) : Latin Modern Mono, inchangé (existe en gras et en italique)
  \setmonofont{lmmono10-regular.otf}[Scale=MatchLowercase,
    BoldFont=lmmonolt10-bold.otf, ItalicFont=lmmono10-italic.otf,
    BoldItalicFont=lmmonolt10-boldoblique.otf]
  \setmathfont{LibertinusMath-Regular.otf}[StylisticSet=8]   % ss08 : intégrales inclinées
}{%
  \PackageWarningNoLine{preamble}{%
    Polices Libertinus introuvables : repli sur Latin Modern.\MessageBreak
    Installer le paquet MiKTeX « libertinus-fonts » pour le rendu prévu}%
  \setmainfont{lmroman10-regular.otf}[
    BoldFont=lmroman10-bold.otf, ItalicFont=lmroman10-italic.otf,
    BoldItalicFont=lmroman10-bolditalic.otf, SmallCapsFont=lmromancaps10-regular.otf]
  \setsansfont{lmsans10-regular.otf}[
    BoldFont=lmsans10-bold.otf, ItalicFont=lmsans10-oblique.otf]
  \setmonofont{lmmono10-regular.otf}[
    BoldFont=lmmonolt10-bold.otf, ItalicFont=lmmono10-italic.otf,
    BoldItalicFont=lmmonolt10-boldoblique.otf]
  \setmathfont{latinmodern-math.otf}
}

%% ---------------------------------------------------------------------
%%  4. Langue : français
%% ---------------------------------------------------------------------
\RequirePackage[french]{babel}
% Listes gérées par enumitem (et non par babel-french) ; l'avertissement
% « legacy list code » est purement informatif, on le désactive.
\frenchsetup{StandardLists=true, SuppressWarning=true,
  CustomiseFigTabCaptions=false,
  SmallCapsFigTabCaptions=false}            % légendes : réglées par caption
\addto\captionsfrench{\def\tablename{Tableau}}   % cohérent avec cleveref
\RequirePackage[autostyle=true]{csquotes}   % \enquote{…} → « … »

%% ---------------------------------------------------------------------
%%  5. Mise en page : A4, marges équilibrées, 11 pt, interligne aéré
%% ---------------------------------------------------------------------
\RequirePackage[a4paper,
  left=2.7cm, right=2.7cm, top=2.9cm, bottom=3.0cm,
  headheight=16pt, headsep=0.9cm, footskip=1.3cm]{geometry}
\linespread{1.05}
\KOMAoptions{parskip=half-}      % paragraphes séparés par un demi-interligne
\RequirePackage[protrusion=true, expansion=true]{microtype}  % justification plus régulière (LuaLaTeX)
\RequirePackage{enumitem}
\setlist{nosep, topsep=0.35em, itemsep=0.2em, leftmargin=*, labelsep=0.6em}
\setlist[itemize,1]{label=\textcolor{AccentColor}{\textendash}}
\setlist[itemize,2]{label=\textcolor{MutedColor}{\textbullet}}
\setlist[enumerate,1]{label=\textcolor{AccentColor}{\arabic*.}}
\setlist[enumerate,2]{label=\textcolor{AccentColor}{(\alph*)}}
\RequirePackage{framed}          % filet vertical de l'environnement solution
\emergencystretch=2em
\hbadness=5000                   % ne signaler que les lignes réellement trop lâches
\vbadness=10000 \vfuzz=3pt       % pages à flottants : pas d'avis sur le remplissage vertical
\raggedbottom                    % bas de page libre : évite les avis « underfull vbox » avec figures et tableaux

%% Titres de sections : Palatino gras, noir (le gras est toujours noir ; le bleu
%% du logo est réservé aux éléments non gras : liens, puces, filet, remarques…)
\setkomafont{disposition}{\sffamily\bfseries\color{TextColor}}
\setkomafont{section}{\Large}
\setkomafont{subsection}{\large}
\setkomafont{subsubsection}{\normalsize}
\setkomafont{paragraph}{\normalfont\normalsize\itshape\color{AccentColor}}  % titres « run-in » en italique romain
\setkomafont{descriptionlabel}{\sffamily\bfseries}
\RedeclareSectionCommand[beforeskip=1.9\baselineskip, afterskip=0.7\baselineskip]{section}
\RedeclareSectionCommand[beforeskip=1.4\baselineskip, afterskip=0.5\baselineskip]{subsection}
\RedeclareSectionCommand[beforeskip=1.1\baselineskip, afterskip=0.3\baselineskip]{subsubsection}
\renewcommand*{\sectionformat}{\thesection\enskip}
\renewcommand*{\subsectionformat}{\thesubsection\enskip}
\renewcommand*{\subsubsectionformat}{\thesubsubsection\enskip}

%% En-têtes et pieds de page discrets
\RequirePackage[headsepline=0.4pt, footsepline=false]{scrlayer-scrpage}
\setkomafont{headsepline}{\color{RuleColor}}
\setkomafont{pageheadfoot}{\normalfont\footnotesize}
\setkomafont{pagenumber}{\normalfont\small\color{TextColor}}
% En-tête sur toutes les pages (sauf la première) : à gauche le titre court en
% italique gris, à droite le cours en romain noir ; filet gris fin dessous.
% Numéro de page centré en pied de page (première page comprise).
\clearpairofpagestyles
\ihead{\itshape\color{MutedColor}\docHeaderLeft}
\ohead{\color{TextColor}\docHeaderRight}
\cfoot*{\pagemark}              % étoile : valable aussi pour la première page (style plain)
\pagestyle{scrheadings}
\addtokomafont{caption}{\small}

%% Notes de bas de page : petit filet gris, marque bleue en exposant, texte 9 pt
\setkomafont{footnote}{\footnotesize\hbadness=10000 }  % pas d'avis « underfull » sur une note d'une ligne
\setkomafont{footnotelabel}{\color{AccentColor}}
\setkomafont{footnotereference}{\color{AccentColor}}
\deffootnote[1.3em]{1.3em}{0pt}{\textsuperscript{\thefootnotemark}\hspace{0.2em}}
\renewcommand*{\footnoterule}{\kern-3pt{\color{RuleColor}\hrule width 0.28\columnwidth height 0.4pt}\kern 2.6pt}
\addtolength{\skip\footins}{4pt}

%% ---------------------------------------------------------------------
%%  6. Réglages mathématiques
%% ---------------------------------------------------------------------
\allowdisplaybreaks[1]
\setlength{\jot}{5pt}
\numberwithin{equation}{section}
\mathtoolsset{mathic=true}
% Virgule décimale en mode math : $0,5$ → 0,5 ; une virgule suivie d'une espace
% reste un séparateur ($a, b$). Les formes tassées $f(x,y)$ restent correctes.
\RequirePackage{icomma}
\setlength{\abovedisplayskip}{9pt plus 3pt minus 4pt}
\setlength{\belowdisplayskip}{9pt plus 3pt minus 4pt}
\setlength{\abovedisplayshortskip}{3pt plus 2pt}
\setlength{\belowdisplayshortskip}{7pt plus 3pt minus 3pt}

%% ---------------------------------------------------------------------
%%  7. Unités et nombres : siunitx v3, virgule décimale française
%% ---------------------------------------------------------------------
\RequirePackage{siunitx}
\sisetup{
  locale                  = FR,
  per-mode                = power,
  inter-unit-product      = \ensuremath{{}\cdot{}},
  range-phrase            = { à~},
  range-units             = single,
  separate-uncertainty    = true,
  uncertainty-mode        = separate,
  exponent-product        = \cdot,
  output-complex-root     = \mathrm{i},
  group-digits            = integer,   % espace fine dans 12 345, pas dans les décimales
  group-minimum-digits    = 4,
}
\DeclareSIUnit{\atmosphere}{atm}
\DeclareSIUnit{\mmHg}{mmHg}
\DeclareSIUnit{\wavenumber}{\per\centi\metre}   % nombre d'onde spectroscopique
\DeclareSIUnit{\gauss}{G}

%% ---------------------------------------------------------------------
%%  8. Figures, tableaux, légendes
%% ---------------------------------------------------------------------
\RequirePackage{graphicx}
\graphicspath{{figures/}}
% Une image ne dépasse jamais la largeur de ligne, n'est jamais déformée,
% et garde sa taille naturelle si elle est plus petite.
\makeatletter
\newcommand*{\maxwidth}{\ifdim\Gin@nat@width>\linewidth\linewidth\else\Gin@nat@width\fi}
\makeatother
\setkeys{Gin}{width=\maxwidth,keepaspectratio}

\RequirePackage{booktabs}
\RequirePackage{array}
\RequirePackage{tabularx}
\RequirePackage{threeparttable}   % notes de tableau : \tnote{a} + environnement tablenotes
\setlength{\heavyrulewidth}{0.8pt}
\setlength{\lightrulewidth}{0.4pt}
\setlength{\cmidrulewidth}{0.3pt}
\renewcommand{\arraystretch}{1.18}
\newcolumntype{Y}{>{\raggedright\arraybackslash}X}   % colonne texte extensible

% caption signale en fin de document tout \captionsetup[type] jamais utilisé
% (ex. [table] dans un document sans tableau). Réglages volontairement
% génériques ici : on vide la liste de contrôle pour éviter ces avis inutiles.
% (Doit être enregistré AVANT le chargement de caption.)
\makeatletter
\AddToHook{enddocument}[tpl-caption-clear]{\@ifundefined{caption@optlist}{}{\gdef\caption@optlist{}}}
\DeclareHookRule{enddocument}{tpl-caption-clear}{before}{caption}
\DeclareHookRule{enddocument}{tpl-caption-clear}{before}{caption3}
\makeatother
\RequirePackage{caption}
\RequirePackage{subcaption}
% Normes éditoriales scientifiques : « Figure 1. Légende… » — étiquette en gras
% (bleu du logo), texte courant, phrase terminée par un point ; figure : légende
% SOUS l'illustration ; tableau et programme : légende AU-DESSUS ; légende d'une
% ligne centrée, de plusieurs lignes justifiée ; corps un cran plus petit.
\captionsetup{
  font           = small,
  labelfont      = {bf,color=TextColor},
  textfont       = {color=TextColor},
  labelsep       = period,
  justification  = justified,
  singlelinecheck= true,
  format         = plain,
  skip           = 8pt,
  margin         = 0.8em,
}
\captionsetup[table]{position=above, skip=6pt}
\captionsetup[figure]{position=below, skip=8pt}
\captionsetup[subfigure]{labelformat=parens, labelfont={bf,color=TextColor}, labelsep=space,
  font=footnotesize, margin=0pt}
\KOMAoptions{captions=tableheading}
% Légende « Programme n » pour les extraits de code (voir \codefile)
\DeclareCaptionType{programme}[Programme][Liste des programmes]
\captionsetup[programme]{position=above, skip=4pt}
% Figures et tableaux : centrés d'office ; texte des tableaux en \small.
% (\@floatboxreset s'exécute DANS la boîte flottante, après \@parboxrestore :
%  c'est le seul endroit où \centering et la taille de police persistent.)
\makeatletter
\AtBeginEnvironment{figure}{\def\@floatboxreset{\reset@font\normalsize\@setminipage\centering}}
\AtBeginEnvironment{table}{\def\@floatboxreset{\reset@font\small\@setminipage\centering}}
\makeatother

%% Graphiques vectoriels (schémas TikZ, courbes pgfplots)
\RequirePackage{tikz}
\usetikzlibrary{arrows.meta,calc,angles,quotes}
\RequirePackage{pgfplots}
\pgfplotsset{compat=newest}
\pgfplotsset{
  every axis/.append style={
    axis line style={color=TextColor, line width=0.4pt},
    tick style={color=TextColor, line width=0.4pt},
    label style={font=\small}, tick label style={font=\footnotesize},
    legend style={font=\footnotesize, draw=RuleColor, fill=none, text=TextColor},
    grid style={color=RuleColor!60, line width=0.3pt},
    /pgf/number format/use comma,          % virgule décimale sur les axes
  },
}

%% Tableaux lus depuis un fichier CSV (résultats d'une simulation Python)
%%   \pgfplotstabletypeset[<options>]{data/resultats.csv}   (voir exemple C)
\RequirePackage{pgfplotstable}
\pgfplotstableset{
  col sep=comma,
  /pgf/number format/use comma,
  every head row/.style={before row=\toprule, after row=\midrule},
  every last row/.style={after row=\bottomrule},
  % notation scientifique « 8,2 × 10⁻¹¹ » : columns/<col>/.style={sci10}
  sci10/.style={/pgf/number format/sci,
    /pgf/number format/sci generic={mantissa sep=\times, exponent={10^{##1}}}},
}

%% Code source (Python, MATLAB, C…) : piton (LuaLaTeX, sans shell-escape, sans minted).
%% Coloration syntaxique par langage : mots-clés, fonctions prédéfinies, chaînes,
%% commentaires, nombres. Les couleurs du code sont indépendantes du thème.
%%   \codefile[options]{Python}{code/x.py}{Légende}{prog:x}   (voir macros.tex)
\RequirePackage{needspace}
\RequirePackage{piton}
\SetPitonStyle{
  Keyword          = \bfseries\color{CodeKeyword},
  Keyword.Governing= \bfseries\color{CodeKeyword},
  Keyword.Constant = \bfseries\color{CodeKeyword},
  Name.Builtin     = \color{CodeBuiltin},
  Name.Function    = \color{TextColor}\bfseries,
  Name.Class       = \color{TextColor}\bfseries,
  Name.Decorator   = \color{CodeBuiltin},
  Name.Namespace   = \color{CodeBuiltin},
  Name.Type        = \color{CodeBuiltin},
  Exception        = \color{CodeString},
  Comment          = \itshape\color{CodeComment},
  String           = \color{CodeString},
  String.Doc       = \itshape\color{CodeComment},
  String.Interpol  = \color{CodeString},
  Number           = \color{CodeNumber},
  Operator         = \color{MutedColor},
  Operator.Word    = \bfseries\color{CodeKeyword},
  Preproc          = \color{CodeBuiltin},
}
\PitonOptions{
  line-numbers,
  line-numbers/format = \tiny\color{CodeLineNo},
  line-numbers/sep    = 8pt,
  line-numbers/absolute,           % avec first-line : numéros du fichier d'origine
  background-color    = FillColor,
  left-margin         = 1.6em,
  break-lines,
  splittable          = 4,         % un long programme peut passer à la page suivante
}
% Langages absents de piton (Python, C, OCaml, SQL le sont d'origine)
\NewPitonLanguage{Matlab}{%
  morekeywords={break,case,catch,classdef,continue,else,elseif,end,for,function,
    global,if,otherwise,parfor,persistent,return,switch,try,while},
  morecomment=[l]\%,
  morestring=[b]',
  morestring=[b]"}
\NewDocumentCommand{\code}{v}{\texttt{\small #1}}   % code en ligne : \code|x = 3|

%% ---------------------------------------------------------------------
%%  9. Bibliographie : biblatex + Biber, style numérique « physique »
%%     (APS/AIP via biblatex-phys, entrées numérotées entre crochets)
%% ---------------------------------------------------------------------
\RequirePackage[
  backend      = biber,
  style        = phys,
  biblabel     = brackets,
  citestyle    = numeric-comp,
  sorting      = none,          % ordre d'apparition des citations
  maxbibnames  = 6, minbibnames = 3,
  doi          = true,
  url          = true,
  isbn         = false,
  eprint       = true,
  giveninits   = true,
  language     = auto,
]{biblatex}
\addbibresource{references.bib}
\DefineBibliographyStrings{french}{bibliography = {Références}}
\defbibheading{bibliography}[\bibname]{\section*{#1}\markright{#1}\addcontentsline{toc}{section}{#1}}
\setlength{\bibitemsep}{0.25\baselineskip}
\urlstyle{same}                 % DOI et URL dans la police du texte
\renewcommand*{\bibfont}{\small}

%% ---------------------------------------------------------------------
%%  10. Hyperliens et références croisées (hyperref puis cleveref)
%% ---------------------------------------------------------------------
\RequirePackage[unicode]{hyperref}
\hypersetup{
  colorlinks   = true,
  linkcolor    = LinkColor,        % gris foncé : cliquable mais sobre
  citecolor    = LinkColor,
  urlcolor     = LinkColor,
  filecolor    = LinkColor,
  bookmarksnumbered = true,
  pdfcreator   = {LuaLaTeX},
}
\RequirePackage[french,nameinlink,noabbrev]{cleveref}
\creflabelformat{equation}{#2\textup{(#1)}#3}

\crefname{equation}{équation}{équations}\Crefname{equation}{Équation}{Équations}
\crefname{figure}{figure}{figures}      \Crefname{figure}{Figure}{Figures}
\crefname{subfigure}{figure}{figures}   \Crefname{subfigure}{Figure}{Figures}
\crefname{table}{tableau}{tableaux}     \Crefname{table}{Tableau}{Tableaux}
\crefname{section}{section}{sections}   \Crefname{section}{Section}{Sections}
\crefname{subsection}{section}{sections}\Crefname{subsection}{Section}{Sections}
\crefname{subsubsection}{section}{sections}\Crefname{subsubsection}{Section}{Sections}
\crefname{definition}{définition}{définitions}     \Crefname{definition}{Définition}{Définitions}
\crefname{theoreme}{théorème}{théorèmes}           \Crefname{theoreme}{Théorème}{Théorèmes}
\crefname{lemme}{lemme}{lemmes}                    \Crefname{lemme}{Lemme}{Lemmes}
\crefname{proposition}{proposition}{propositions}  \Crefname{proposition}{Proposition}{Propositions}
\crefname{corollaire}{corollaire}{corollaires}     \Crefname{corollaire}{Corollaire}{Corollaires}
\crefname{remarque}{remarque}{remarques}           \Crefname{remarque}{Remarque}{Remarques}
\crefname{exemple}{exemple}{exemples}              \Crefname{exemple}{Exemple}{Exemples}
\crefname{exercice}{exercice}{exercices}           \Crefname{exercice}{Exercice}{Exercices}
\crefname{question}{question}{questions}           \Crefname{question}{Question}{Questions}
\crefname{enumi}{point}{points}                    \Crefname{enumi}{Point}{Points}
\crefname{programme}{programme}{programmes}      \Crefname{programme}{Programme}{Programmes}

%% ---------------------------------------------------------------------
%%  11. Informations du document et titre (voir style/titre.tex)
%% ---------------------------------------------------------------------
% !TeX encoding = UTF-8
%% =====================================================================
%%  style/titre.tex — Informations du document (\docsetup) et mise en page
%%  du titre : bandeau [logo | institution | type de document], titre
%%  compact (par défaut) ou page de titre complète.
%%  Chargé automatiquement par style/preamble.tex. À ne pas modifier.
%%
%%  \docsetup{clé = valeur, …}
%%    - identité permanente  → config.tex  (une seule fois)
%%    - infos du document    → début du fichier .tex de chaque travail
%%  Clés : type titlestyle logo brouillon etudiant universite faculte
%%         departement filiere annee cours enseignant titre titrecourt
%%         soustitre date numero binome groupe seance typelabel
%% =====================================================================

\pgfkeys{
  /doc/.is family, /doc,
  type/.store in             = \docType,
  titlestyle/.store in       = \docTitleStyle,
  logo/.store in             = \docLogo,
  brouillon/.store in        = \docDraft,
  etudiant/.store in         = \docEtudiant,
  universite/.store in       = \docUniversite,
  faculte/.store in          = \docFaculte,
  departement/.store in      = \docDepartement,
  filiere/.store in          = \docFiliere,
  annee/.store in            = \docAnnee,
  cours/.store in            = \docCours,
  enseignant/.store in       = \docEnseignant,
  titre/.store in            = \docTitre,
  titrecourt/.store in       = \docTitreCourt,
  soustitre/.store in        = \docSousTitre,
  date/.store in             = \docDate,
  numero/.store in           = \docNumero,
  binome/.store in           = \docBinome,
  groupe/.store in           = \docGroupe,
  seance/.store in           = \docSeance,
  typelabel/.store in        = \docTypeLabel,
}

% \docsetup accepte lignes vides et commentaires : les \par éventuels sont
% retirés avant d'être transmis à pgfkeys (qui, lui, ne les supporte pas).
\ExplSyntaxOn
\newcommand{\docsetup}[1]{%
  \tl_set:Nn \l_tmpa_tl {#1}
  \tl_remove_all:Nn \l_tmpa_tl {\par}
  \tl_put_left:Nn \l_tmpa_tl {/doc/.cd,}
  \exp_args:NV \pgfkeys \l_tmpa_tl
}
\ExplSyntaxOff

% Valeurs par défaut (toute clé non renseignée est simplement omise)
\docsetup{type=assignment, titlestyle=block, logo={}, brouillon=false,
  etudiant={}, universite={}, faculte={}, departement={}, filiere={}, annee={},
  cours={}, enseignant={}, titre={Titre du document}, titrecourt={},
  soustitre={}, date={\today}, numero={}, binome={}, groupe={}, seance={},
  typelabel={}}

%% ---------------------------------------------------------------------
%%  Outils internes
%% ---------------------------------------------------------------------
\newcommand*{\ifdoctype}[1]{\ifdefstring{\docType}{#1}}

% Intitulé du type, volontairement COURT (modifiable avec la clé typelabel)
\newcommand*{\docTypeName}{%
  \ifdefempty{\docTypeLabel}
    {\ifdoctype{tp}{Compte rendu}
      {\ifdoctype{simulation}{Compte rendu}
        {\ifdoctype{rapport}{Rapport}{Devoir à la maison}}}}
    {\docTypeLabel}}

\newcommand*{\docHeaderLeft}{\ifdefempty{\docTitreCourt}{\docTitre}{\docTitreCourt}}
\newcommand*{\docHeaderRight}{\docCours}

% Logo facultatif : fichier introuvable → avertissement, titre sans logo
\newbool{doclogo}
\newlength{\doclogow}
\setlength{\doclogow}{2.7cm}
\newcommand*{\docCheckLogo}{%
  \ifdefempty{\docLogo}{\boolfalse{doclogo}}{%
    \IfFileExists{\docLogo}{\booltrue{doclogo}}{%
      \boolfalse{doclogo}%
      \PackageWarning{titre}{Logo introuvable : \docLogo}}}}

% Élément « étiquette valeur » de la ligne d'informations (omis si vide)
\newcounter{docmetaN}
% Séparateur « | » gris entre deux éléments de la ligne d'informations. Il n'est
% coupable d'aucune coupure disgracieuse : si la ligne passe à la ligne suivante
% à cet endroit, la barre disparaît (discrétionnaire) au lieu de rester isolée.
\newsavebox{\docbarbox}
\AtBeginDocument{\sbox{\docbarbox}{\textcolor{RuleColor}{|}}}
\newcommand*{\docsep}{\nobreak\hspace{0.8em}\discretionary{}{}{\copy\docbarbox}\nobreak\hspace{0.8em}}
\makeatletter
\newcommand*{\docitem}[2]{%
  \protected@xdef\docmeta@val{#2}%
  \ifdefempty{\docmeta@val}{}{%
    \ifnum\value{docmetaN}>0 \docsep\fi
    \stepcounter{docmetaN}%
    {\sffamily\footnotesize\color{MutedColor}#1}\hspace{0.45em}\docmeta@val}}
% même chose en tableau (page de titre complète)
\newcommand*{\docrow}[2]{%
  \protected@xdef\docmeta@val{#2}%
  \ifdefempty{\docmeta@val}{}{%
    {\sffamily\footnotesize\scshape\color{MutedColor}#1}&\docmeta@val\\}}
\makeatother

% Étiquette du numéro : « TP n° », « Simulation n° », « N° »
\newcommand*{\docNumeroLabel}{%
  \ifdoctype{tp}{TP n\textsuperscript{o}}
    {\ifdoctype{simulation}{Simulation n\textsuperscript{o}}{N\textsuperscript{o}}}}

\newcommand*{\docmetaline}{%
  \setcounter{docmetaN}{0}%
  \docitem{\docNumeroLabel}{\docNumero}%
  \docitem{Étudiant\ifdoctype{tp}{(e)}{}}{\docEtudiant}%
  \ifdoctype{tp}{\docitem{Binôme}{\docBinome}\docitem{Groupe}{\docGroupe}}{}%
  \docitem{Enseignant}{\docEnseignant}%
  \ifdoctype{tp}{\docitem{Séance}{\docSeance}}{}%
  \docitem{Date}{\docDate}}

\newcommand*{\docmetatable}{%
  \begingroup\renewcommand{\arraystretch}{1.3}%
  \begin{tabular}{@{}r@{\hspace{1.1em}}l@{}}
    \docrow{\docNumeroLabel}{\docNumero}%
    \docrow{Étudiant\ifdoctype{tp}{(e)}{}}{\docEtudiant}%
    \ifdoctype{tp}{\docrow{Binôme}{\docBinome}\docrow{Groupe}{\docGroupe}}{}%
    \docrow{Enseignant}{\docEnseignant}%
    \ifdoctype{tp}{\docrow{Séance}{\docSeance}}{}%
    \docrow{Date}{\docDate}%
  \end{tabular}\endgroup}

%% ---------------------------------------------------------------------
%%  Bandeau : [LOGO à gauche] [bloc gauche]            [bloc droit]
%%  Les deux blocs de texte sont construits à l'identique (3 lignes de
%%  même corps, mêmes interlignes) et centrés verticalement sur le logo :
%%  leurs lignes sont donc exactement au même niveau, à gauche (texte
%%  aligné à gauche) et à droite (texte aligné à droite).
%% ---------------------------------------------------------------------
\newcommand*{\docFacDept}{%
  \docFaculte\ifdefempty{\docDepartement}{}{\ifdefempty{\docFaculte}{}{\ \textperiodcentered\ }\docDepartement}}

\newcommand*{\docLeftBlock}{%
  \begingroup\renewcommand{\arraystretch}{1.1}%
  \begin{tabular}[c]{@{}l@{}}
    {\sffamily\small\bfseries\scshape\color{TextColor}\docUniversite\strut}\\
    {\sffamily\footnotesize\color{MutedColor}\docFacDept\strut}\\
    {\sffamily\footnotesize\color{MutedColor}\docFiliere\strut}
  \end{tabular}\endgroup}

\newcommand*{\docRightBlock}{%
  \begingroup\renewcommand{\arraystretch}{1.1}%
  \begin{tabular}[c]{@{}r@{}}
    {\sffamily\small\bfseries\scshape\color{TextColor}\docTypeName\strut}\\
    {\sffamily\footnotesize\color{MutedColor}\docCours\strut}\\
    {\sffamily\footnotesize\color{MutedColor}\docAnnee\strut}
  \end{tabular}\endgroup}

\newcommand*{\docHeaderBand}{%
  \docCheckLogo
  \noindent
  \ifbool{doclogo}
    {\parbox[c]{\doclogow}{\includegraphics[width=\doclogow]{\docLogo}}\hspace{1.2em}}{}%
  \docLeftBlock\hfill\docRightBlock\par}

% Filet bleu marine (couleur du logo) sous le bandeau.
\newcommand*{\docTopRule}{%
  {\color{AccentColor}\rule{\linewidth}{0.8pt}}\par}

%% ---------------------------------------------------------------------
%%  Bloc de titre compact : bandeau, filet, titre, sous-titre, informations
%% ---------------------------------------------------------------------
\newcommand*{\docTitleBlock}{%
  \thispagestyle{plain.scrheadings}%
  \begingroup
  \docHeaderBand
  \vspace{0.7em}\docTopRule\vspace{0.9em}
  \raggedright
  {\fontsize{20}{24}\selectfont\bfseries\color{TextColor}\docTitre\par}
  \ifdefempty{\docSousTitre}{}{\vspace{0.25em}{\large\itshape\color{MutedColor}\docSousTitre\par}}
  \vspace{0.7em}
  {\docmetaline\par}
  \vspace{0.5em}{\color{RuleColor}\rule{\linewidth}{0.4pt}}\par\vspace{0.9em}
  \endgroup}

%% ---------------------------------------------------------------------
%%  Page de titre complète (titlestyle = page)
%% ---------------------------------------------------------------------
\newcommand*{\docTitlePage}{%
  \begin{titlepage}\thispagestyle{empty}\raggedright
  \setlength{\doclogow}{3.4cm}%
  \docHeaderBand
  \vspace{0.8em}\docTopRule
  \vfill
  {\sffamily\small\bfseries\scshape\color{TextColor}\docTypeName%
   \ifdefempty{\docCours}{}{\ \textperiodcentered\ \docCours}\par}\vspace{0.8em}
  {\fontsize{30}{35}\selectfont\bfseries\color{TextColor}\docTitre\par}
  \ifdefempty{\docSousTitre}{}{\vspace{0.6em}{\Large\itshape\color{MutedColor}\docSousTitre\par}}
  \vspace{1.2em}{\color{RuleColor}\rule{0.3\linewidth}{0.8pt}}
  \vfill
  \docmetatable
  \ifdefempty{\docAnnee}{}{\par\vspace{0.4em}{\sffamily\footnotesize\color{MutedColor}\docAnnee}}
  \par\vspace{1.2em}{\color{AccentColor}\rule{\linewidth}{0.8pt}}
  \end{titlepage}}

% À placer juste après \begin{document}
\newcommand*{\MakeDocTitle}{%
  \ifdefstring{\docTitleStyle}{page}{\docTitlePage}{\docTitleBlock}}

%% ---------------------------------------------------------------------
%%  Filigrane « BROUILLON » (clé brouillon = true) et métadonnées PDF
%% ---------------------------------------------------------------------
\AtBeginDocument{%
  \hypersetup{pdftitle={\docTitre}, pdfauthor={\docEtudiant}, pdfsubject={\docCours}}%
  \ifdefstring{\docDraft}{true}{%
    \AddToHook{shipout/background}{%
      \put(0.5\paperwidth,-0.5\paperheight){\makebox(0,0){%
        \rotatebox{52}{\fontsize{96}{96}\selectfont\sffamily\bfseries
          \color[gray]{0.9}BROUILLON}}}}}{}}


% !TeX encoding = UTF-8
%% =====================================================================
%%  style/macros.tex — Commandes et environnements personnels.
%%
%%  RÈGLE : toutes les commandes sont créées avec \NewDocumentCommand,
%%  \DeclarePairedDelimiter, \DeclareMathOperator ou \newtheorem, qui
%%  produisent une ERREUR si le nom existe déjà. Aucune commande LaTeX,
%%  AMS ou unicode-math n'est donc redéfinie en silence.
%%  physics / physics2 ne sont PAS chargés (ils redéfinissent de nombreux
%%  noms usuels) : les commandes utiles sont fournies explicitement ici.
%%
%%  SOMMAIRE
%%    1. Notations mathématiques et physiques
%%    2. Environnements de type théorème (définition, théorème, …)
%%    3. Démonstration et fin de démonstration
%%    4. Exercices, questions, solutions
%%    5. Blocs de compte rendu de TP (objectif, principe, protocole, …)
%%    6. Divers (résumé, incertitudes)
%% =====================================================================

%% ---------------------------------------------------------------------
%%  1. NOTATIONS MATHÉMATIQUES ET PHYSIQUES
%%     Convention typographique (proche de ISO 80000-2) :
%%       grandeurs scalaires ........ italique            a, E, \omega
%%       vecteurs ................... gras droit         \vect{v}
%%       vecteurs unitaires ......... gras droit + ^     \uvect{e}
%%       matrices ................... gras droit          \mat{A}
%%       tenseurs ................... gras sans empatt.   \tens{T}
%%       opérateurs quantiques ...... italique + chapeau  \op{H}
%%       constantes i, e, d ......... droit               \ii  \ee  \dd
%% ---------------------------------------------------------------------

% --- Constantes mathématiques droites -------------------------------
\NewDocumentCommand{\ii}{}{\symup{i}}       % unité imaginaire
\NewDocumentCommand{\ee}{}{\symup{e}}       % base de l'exponentielle : \ee^{-x}
\NewDocumentCommand{\dd}{}{\symup{d}}       % différentielle : \int f \dd x

% --- Ensembles de nombres -------------------------------------------
\NewDocumentCommand{\RR}{}{\symbb{R}}
\NewDocumentCommand{\CC}{}{\symbb{C}}
\NewDocumentCommand{\NN}{}{\symbb{N}}
\NewDocumentCommand{\ZZ}{}{\symbb{Z}}
\NewDocumentCommand{\QQ}{}{\symbb{Q}}

% --- Vecteurs, matrices, tenseurs, opérateurs -----------------------
\NewDocumentCommand{\vect}{m}{\symbf{#1}}                  % \vect{v}, \vect{\omega}
\NewDocumentCommand{\uvect}{m}{\hat{\symbf{#1}}}           % \uvect{n}
\NewDocumentCommand{\ex}{}{\uvect{e}_{x}}                 % vecteurs de base
\NewDocumentCommand{\ey}{}{\uvect{e}_{y}}
\NewDocumentCommand{\ez}{}{\uvect{e}_{z}}
\NewDocumentCommand{\er}{}{\uvect{e}_{r}}
\NewDocumentCommand{\etheta}{}{\uvect{e}_{\theta}}
\NewDocumentCommand{\ephi}{}{\uvect{e}_{\varphi}}
\NewDocumentCommand{\mat}{m}{\symbfup{#1}}                % \mat{A}
\NewDocumentCommand{\tens}{m}{\symbfsfit{#1}}             % \tens{T}_{ij}
\NewDocumentCommand{\op}{m}{\hat{#1}}                     % \op{H}, \op{x}
\NewDocumentCommand{\adj}{}{^{\dagger}}                   % adjoint : \op{a}\adj
\NewDocumentCommand{\cc}{}{^{*}}                          % conjugué complexe : z\cc
\NewDocumentCommand{\hc}{}{\text{h.c.}}                   % « hermitien conjugué »
\NewDocumentCommand{\transpose}{}{^{\mathsf{T}}}          % transposée : \mat{A}\transpose

% --- Fonctions et opérateurs usuels ---------------------------------
\DeclareMathOperator{\Tr}{Tr}
\DeclareMathOperator{\sgn}{sgn}
\DeclareMathOperator{\diag}{diag}
\DeclareMathOperator{\Reel}{Re}                           % \Reel z (partie réelle)
\DeclareMathOperator{\Imag}{Im}                           % \Imag z (partie imaginaire)
\NewDocumentCommand{\bigO}{m}{\symcal{O}{\left(#1\right)}}           % \bigO{x^2}

% --- Dérivées -------------------------------------------------------
%   \dv{f}{x}  \dv[2]{f}{x}            dérivée ordinaire d'ordre n (défaut 1)
%   \pdv{f}{x} \pdv[2]{f}{x}           dérivée partielle d'ordre n
%   \pdv{f}{x}{y}                      dérivée partielle croisée
%   \fdv{S}{q}                         dérivée fonctionnelle
\NewDocumentCommand{\dv}{O{1} m m}{%
  \frac{\dd\ifnum#1>1 ^{#1}\fi #2}{\dd #3\ifnum#1>1 ^{#1}\fi}}
\NewDocumentCommand{\pdv}{O{1} m m g}{%
  \IfNoValueTF{#4}
    {\frac{\partial\ifnum#1>1 ^{#1}\fi #2}{\partial #3\ifnum#1>1 ^{#1}\fi}}
    {\frac{\partial^{2} #2}{\partial #3\,\partial #4}}}
\NewDocumentCommand{\fdv}{m m}{\frac{\delta #1}{\delta #2}}

% --- Opérateurs différentiels (opérateur nabla en gras) -------------
%   \grad f     \diverg\vect{E}     \curl\vect{B}     \lapl f
\NewDocumentCommand{\grad}{}{\symbf{\nabla}}
\NewDocumentCommand{\diverg}{}{\symbf{\nabla}\cdot}
\NewDocumentCommand{\curl}{}{\symbf{\nabla}\times}
\NewDocumentCommand{\lapl}{}{\nabla^{2}}

% --- Produits, normes, valeurs absolues -----------------------------
%   Toutes les paires de délimiteurs acceptent une taille explicite
%   \abs[\big]{x}, ou la taille automatique \abs*{\frac{a}{b}}.
\NewDocumentCommand{\dotp}{}{\mathbin{\cdot}}            % \vect{a}\dotp\vect{b}
\NewDocumentCommand{\crossp}{}{\mathbin{\times}}         % \vect{a}\crossp\vect{b}
\DeclarePairedDelimiter{\abs}{\lvert}{\rvert}
\DeclarePairedDelimiter{\norm}{\lVert}{\rVert}
\DeclarePairedDelimiterX{\pscal}[2]{\langle}{\rangle}{#1,#2}   % produit scalaire ⟨a,b⟩

% --- Commutateurs, anticommutateurs, crochets de Poisson ------------
\DeclarePairedDelimiterX{\comm}[2]{\lbrack}{\rbrack}{#1,#2}    % [A,B]
\DeclarePairedDelimiterX{\acomm}[2]{\lbrace}{\rbrace}{#1,#2}   % {A,B}
\DeclarePairedDelimiterX{\poisson}[2]{\lbrace}{\rbrace}{#1,#2}  % {f,g}_P (usage mécanique)

% --- Notation de Dirac, valeurs moyennes, éléments de matrice -------
\DeclarePairedDelimiter{\ket}{\lvert}{\rangle}                         % |ψ⟩
\DeclarePairedDelimiter{\bra}{\langle}{\rvert}                         % ⟨ψ|
\DeclarePairedDelimiterX{\braket}[2]{\langle}{\rangle}{#1\delimsize\vert#2}   % ⟨φ|ψ⟩
\DeclarePairedDelimiterX{\ketbra}[2]{\lvert}{\rvert}{#1\delimsize\rangle\delimsize\langle#2} % |φ⟩⟨ψ|
\DeclarePairedDelimiterX{\mel}[3]{\langle}{\rangle}{#1\delimsize\vert#2\delimsize\vert#3}    % ⟨m|A|n⟩
\DeclarePairedDelimiter{\mean}{\langle}{\rangle}                       % ⟨x⟩ (moyenne quelconque)
\DeclarePairedDelimiterX{\evstate}[2]{\langle}{\rangle}{#2\delimsize\vert#1\delimsize\vert#2}
%   \ev{A}       → ⟨A⟩            (valeur moyenne)
%   \ev{A}{\psi} → ⟨ψ|A|ψ⟩       (espérance de A dans l'état ψ)
%   \ev*{…}{…}  → délimiteurs de taille automatique
\NewDocumentCommand{\ev}{s m g}{%
  \IfNoValueTF{#3}
    {\IfBooleanTF{#1}{\mean*{#2}}{\mean{#2}}}
    {\IfBooleanTF{#1}{\evstate*{#2}{#3}}{\evstate{#2}{#3}}}}

% --- Incertitudes de mesure -----------------------------------------
\NewDocumentCommand{\incert}{m}{u\mathopen{}\left(#1\right)}         % u(g)
\NewDocumentCommand{\incertrel}{m}{\frac{u\mathopen{}\left(#1\right)}{#1}}  % u(g)/g
\NewDocumentCommand{\ecartnorm}{m m}{%                                % écart normalisé
  \frac{\abs{#1-#2}}{\sqrt{u^{2}\mathopen{}\left(#1\right)+u^{2}\mathopen{}\left(#2\right)}}}
%   \qtyerr{9.81}{0.05}{\metre\per\second\squared} → (9,81 ± 0,05) m·s⁻²
\NewDocumentCommand{\qtyerr}{m m m}{\qty{#1 +- #2}{#3}}

%% ---------------------------------------------------------------------
%%  2. ENVIRONNEMENTS DE TYPE THÉORÈME (sobres : pas de boîtes)
%%     Titre en sans empattement bleu marine, filets absents.
%%     Chaque environnement accepte un titre optionnel :
%%        \begin{theoreme}[Inégalité de Cauchy–Schwarz] … \end{theoreme}
%%     Définition, théorème, lemme, proposition, corollaire partagent un
%%     même compteur (numérotation par section : 1.1, 1.2, …).
%% ---------------------------------------------------------------------
\newtheoremstyle{acadplain}                 % énoncés en italique
  {0.9\baselineskip plus 2pt minus 2pt}{0.7\baselineskip plus 2pt minus 2pt}
  {\itshape}{0pt}
  {\sffamily\bfseries\color{TextColor}}{.}
  {0.6em}
  {\thmname{#1}\thmnumber{ #2}\thmnote{ {\normalfont\color{TextColor}(#3)}}}

\newtheoremstyle{acaddef}                   % énoncés en romain
  {0.9\baselineskip plus 2pt minus 2pt}{0.7\baselineskip plus 2pt minus 2pt}
  {\normalfont}{0pt}
  {\sffamily\bfseries\color{TextColor}}{.}
  {0.6em}
  {\thmname{#1}\thmnumber{ #2}\thmnote{ {\normalfont\color{TextColor}(#3)}}}

\newtheoremstyle{acadremark}                % remarques, exemples : titre en italique
  {0.7\baselineskip plus 2pt minus 2pt}{0.6\baselineskip plus 2pt minus 2pt}
  {\normalfont}{0pt}
  {\itshape\color{AccentColor}}{.}
  {0.6em}
  {\thmname{#1}\thmnumber{ #2}\thmnote{ {\upshape\color{TextColor}(#3)}}}

\theoremstyle{acadplain}
\newtheorem{theoreme}{Théorème}[section]
\newtheorem{lemme}[theoreme]{Lemme}
\newtheorem{proposition}[theoreme]{Proposition}
\newtheorem{corollaire}[theoreme]{Corollaire}
\theoremstyle{acaddef}
\newtheorem{definition}[theoreme]{Définition}
\theoremstyle{acadremark}
\newtheorem{remarque}{Remarque}[section]
\newtheorem{exemple}{Exemple}[section]
\newtheorem*{note}{Note}                    % remarque non numérotée (ex. avant la section 1)

%% ---------------------------------------------------------------------
%%  3. DÉMONSTRATION ET SYMBOLE DE FIN
%%     \begin{demonstration} … \end{demonstration}  (alias de « proof »)
%%     Dans une équation finale : \qedhere
%% ---------------------------------------------------------------------
\addto\captionsfrench{\def\proofname{Démonstration}}
\renewcommand{\qedsymbol}{%
  {\color{MutedColor}\fboxsep=0pt \fboxrule=0.4pt
   \fbox{\rule{0pt}{0.62em}\rule{0.62em}{0pt}}}}
\makeatletter
\renewenvironment{proof}[1][\proofname]{%
  \par\pushQED{\qed}\normalfont
  \topsep=0.7\baselineskip plus 2pt minus 2pt \relax
  \trivlist\item[\hskip\labelsep{\itshape\color{AccentColor}#1.}]\ignorespaces
}{\popQED\endtrivlist\@endpefalse}
\makeatother
\let\demonstration\proof
\let\enddemonstration\endproof

%% ---------------------------------------------------------------------
%%  4. EXERCICES, QUESTIONS, SOLUTIONS
%%     L'énoncé (exercice/question) est composé en noir, avec un titre
%%     bleu marine ; la solution est repérable à son filet gris vertical.
%%        \begin{exercice}[Titre] … \end{exercice}
%%        \begin{question}[points] … \end{question}       (Question 1.1, 1.2…)
%%        \begin{solution} … \end{solution}
%% ---------------------------------------------------------------------
\newtheoremstyle{acadexo}
  {1.3\baselineskip plus 3pt minus 2pt}{0.5\baselineskip plus 2pt minus 2pt}
  {\normalfont}{0pt}
  {\sffamily\bfseries\large\color{TextColor}}{.}
  {0.5em}
  {\thmname{#1}\thmnumber{ #2}\thmnote{ \textemdash\ {\normalfont\bfseries\color{TextColor}#3}}}
\newtheoremstyle{acadquestion}
  {0.9\baselineskip plus 2pt minus 2pt}{0.4\baselineskip plus 2pt minus 2pt}
  {\normalfont}{0pt}
  {\sffamily\bfseries\color{TextColor}}{.}
  {0.6em}
  {\thmname{#1}\thmnumber{ #2}\thmnote{ {\normalfont\sffamily\footnotesize\color{MutedColor}[#3]}}}
\theoremstyle{acadexo}
\newtheorem{exercice}{Exercice}
\theoremstyle{acadquestion}
\newtheorem{question}{Question}[exercice]

\NewDocumentEnvironment{solution}{O{Solution}}{%
  \par\addvspace{0.2\baselineskip}%
  \def\FrameCommand{\hspace{0.3em}{\color{RuleColor}\vrule width 0.9pt}\hspace{0.9em}}%
  \MakeFramed{\advance\hsize-\width \FrameRestore}%
  \noindent{\itshape\color{TextColor}#1.}\enspace\ignorespaces
}{%
  \par\endMakeFramed\par\addvspace{0.4\baselineskip}}

%% ---------------------------------------------------------------------
%%  5. BLOCS DE COMPTE RENDU DE TP
%%     Chaque bloc ouvre une VRAIE section numérotée (présente dans la
%%     table des matières, référençable avec \label juste après \begin).
%%        \begin{objectif} … \end{objectif}
%%        \begin{objectif}[Autre titre] … \end{objectif}   titre modifié
%%        \begin{objectif*} … \end{objectif*}              section non numérotée
%% ---------------------------------------------------------------------
\NewDocumentCommand{\DeclareTPBlock}{m m}{%
  \NewDocumentEnvironment{#1}{o}{\section{\IfNoValueTF{##1}{#2}{##1}}}{\par}%
  \NewDocumentEnvironment{#1*}{o}{\section*{\IfNoValueTF{##1}{#2}{##1}}}{\par}}
\DeclareTPBlock{objectif}{Objectif}
\DeclareTPBlock{materiel}{Matériel}
\DeclareTPBlock{principe}{Principe théorique}
\DeclareTPBlock{protocole}{Protocole expérimental}
\DeclareTPBlock{resultats}{Résultats expérimentaux}
\DeclareTPBlock{analyse}{Analyse et discussion}
\DeclareTPBlock{conclusion}{Conclusion}
% Blocs pour un compte rendu de SIMULATION NUMÉRIQUE (type = simulation)
\DeclareTPBlock{modele}{Modèle physique}
\DeclareTPBlock{methode}{Méthode numérique}
\DeclareTPBlock{implementation}{Implémentation}
\DeclareTPBlock{validation}{Validation et convergence}

%% ---------------------------------------------------------------------
%%  5 bis. CODE SOURCE (coloration syntaxique par langage, via piton)
%%     \codefile[options]{Langage}{fichier}{Légende}{étiquette}
%%       Langage : Python | C | Matlab | OCaml | SQL (voir style/preamble.tex)
%%       options : first-line=22, last-line=47 (extrait), line-numbers=false…
%%       Référence : \cref{étiquette} → « programme 1 »
%%     Code en ligne : \code|x = 3|    Code écrit à la main : environnement Piton.
%% ---------------------------------------------------------------------
\NewDocumentCommand{\codefile}{O{} m m m m}{%
  \par\addvspace{0.8\baselineskip}\begingroup\needspace{6\baselineskip}%
  \captionof{programme}{#4}\label{#5}\nopagebreak\vspace{0.3em}%
  \footnotesize\PitonInputFile[language=#2,#1]{#3}%
  \endgroup\par\addvspace{0.4\baselineskip}}

%% ---------------------------------------------------------------------
%%  6. DIVERS
%% ---------------------------------------------------------------------
% Résumé en tête de document, sous le bloc de titre
\NewDocumentEnvironment{resume}{}{%
  \begin{quote}\small
  {\sffamily\bfseries\scshape\color{TextColor}Résumé.}\enspace\ignorespaces
}{\end{quote}}

%% ---------------------------------------------------------------------
%%  7. SOURCES ET NOTES
%%     \source{Texte}  : ligne de crédit sous une figure ou un tableau
%%                       (norme éditoriale : « Source : … » ou « Adapté de [3] »)
%%     Notes de bas de page : \footnote{…}
%%     Notes de tableau (threeparttable) :
%%        \begin{table} \caption{…}
%%          \begin{threeparttable}
%%            \begin{tabular}{…} … {$u(T)$\tnote{a}} … \end{tabular}
%%            \begin{tablenotes}[flushleft]\footnotesize \item[a] Texte de la note. \end{tablenotes}
%%          \end{threeparttable}
%%        \end{table}
%% ---------------------------------------------------------------------
\NewDocumentCommand{\source}{m}{%
  \par\vspace{0.3em}\noindent{\footnotesize\itshape Source : #1\par}}

\geometry{left=1.2cm,right=1.2cm,top=0.5cm,bottom=0.6cm,headheight=15pt,headsep=0pt,footskip=0pt}
\raggedbottom
\pagestyle{empty}
\setlength{\parindent}{0pt}
\linespread{1}
\emergencystretch=3em
\setlength{\tabcolsep}{1pt}

\newcommand{\cheatsize}{\fontsize{8}{9.4}\selectfont}
\newcommand{\sect}[1]{\par\vspace{0.5em}{\sffamily\bfseries\small #1}\par\vspace{-0.2em}{\color{RuleColor}\rule{\linewidth}{0.4pt}}\par\vspace{0.1em}}
\newcommand{\cmd}{\cheatsize\verb}                 % code littéral : \cmd~\vect{v}~
\newcommand{\w}[1]{\texttt{\cheatsize #1}}         % mots simples (clés, noms d'environnements)
\newcolumntype{P}[1]{>{\raggedright\arraybackslash}p{#1}}
\newenvironment{tbl}[1]{\cheatsize\renewcommand{\arraystretch}{1.1}\begin{tabular}{@{}#1@{}}}{\end{tabular}\par}

\begin{document}
\nocite{griffiths2018}\raggedright   % évite l'avertissement Biber « aucune citation »
{\sffamily\bfseries\Large Aide-mémoire}\hfill
{\sffamily\small Template physique (LuaLaTeX + Biber) — compiler : \textbf{F5} ou \texttt{compiler.bat}}\par
\vspace{-0.2em}{\color{TextColor}\rule{\linewidth}{0.8pt}}\par\vspace{0.1em}

%% ============================ COLONNE GAUCHE ============================
\begin{minipage}[t]{0.465\textwidth}
\cheatsize

\sect{1. Nouveau travail}
\begin{tbl}{P{8.4cm}}
Copier \cmd~main.tex~, le renommer, régler \cmd~\docsetup{…}~ en haut du fichier.\\
Identité et logo : \cmd~config.tex~ (une seule fois).\\
Compiler : \textbf{F5} (TeXstudio) ou double-clic \cmd~compiler.bat~.\\
Simulation : \cmd~python code/simulation.py~ puis compiler.\\
\end{tbl}

\sect{2. Clés de \texttt{\textbackslash docsetup}}
\begin{tbl}{P{2.9cm}P{5.5cm}}
\w{type} & \w{assignment}, \w{tp}, \w{simulation}, \w{rapport}\\
\w{titre soustitre} & titre et sous-titre ; \w{titrecourt} : en-tête de page\\
\w{cours enseignant} & module, enseignant\\
\w{date} & \cmd~{\today}~ par défaut\\
\w{numero} & numéro du TP / de la simulation\\
\w{binome groupe seance} & pour le type \w{tp}\\
\w{titlestyle} & \w{block} (compact) ou \w{page}\\
\w{brouillon} & \w{true} : filigrane BROUILLON\\
\w{logo} & chemin de l'image (vide = aucun)\\
\w{etudiant universite faculte departement filiere annee} & identité, dans \w{config.tex}\\
\w{typelabel} & impose l'intitulé du type\\
\end{tbl}

\sect{3. Blocs de section (numérotés)}
\begin{tbl}{P{2.0cm}P{6.4cm}}
TP & \w{objectif materiel principe protocole resultats analyse conclusion}\\
Simulation & \w{objectif modele methode implementation resultats validation analyse conclusion}\\
Variantes & \cmd~\begin{principe}[Titre]~ titre libre ;\ \cmd~\begin{objectif*}~ non numéroté\\
\end{tbl}

\sect{4. Environnements}
\begin{tbl}{P{2.0cm}P{6.4cm}}
Énoncés & \w{definition theoreme lemme proposition corollaire} ; titre : \cmd~[Titre]~\\
Commentaires & \w{remarque exemple} (numérotés), \w{note} (non numérotée)\\
Preuve & \w{demonstration} ; fin automatique, ou \cmd~\qedhere~\\
Devoir & \w{exercice} \cmd~[Titre]~, \w{question} \cmd~[3 pts]~, \w{solution} \cmd~[Réponse]~\\
Résumé & \w{resume}\\
\end{tbl}

\sect{5. Références croisées et bibliographie}
\begin{tbl}{P{3.0cm}P{5.4cm}}
\cmd~\label{eq:x}~ & juste \emph{après} \cmd~\begin{…}~ ou \cmd~\caption~\\
\cmd~\cref{eq:x}~ & « équation (1.2) » ; \cmd~\Cref~ en début de phrase\\
\cmd~\eqref{eq:x}~ & « (1.2) » seul\\
\cmd~\cref{a,b}~ & « équations (1.2) et (1.3) »\\
\cmd~\cite{cle}~ & [1] ; entrée dans \w{references.bib}\\
\end{tbl}

\sect{6. Figures, tableaux, code}
\begin{tbl}{P{1.4cm}P{7.0cm}}
Image & \cmd~\includegraphics{courbe}~ (dossier \w{figures/})\\
Figure & \cmd~\begin{figure}[htbp]~, \cmd~\caption~, \cmd~\label~\\
Sous-fig. & \cmd~\begin{subfigure}[b]{0.48\textwidth}~\\
Tableau & colonnes \cmd~S[table-format=3.1]~, lignes \cmd~\toprule~ \cmd~\midrule~ \cmd~\bottomrule~\\
CSV & \cmd~\pgfplotstabletypeset{f.csv}~ ; options \w{fixed, fixed zerofill}\\
Courbe & \cmd~\addplot table[x=a,y=b]{f.csv}~\\
Code & \cmd~\codefile{Python}{f.py}{Légende}{prog:x}~ ; langages \w{Python C Matlab OCaml SQL} ; extrait : \cmd~[first-line=22,last-line=47]~\\
En ligne & \cmd~\code|x = 3|~ ; \cmd~\cref{prog:x}~ : « programme 1 »\\
Solution & figure : \cmd~\captionof{figure}{…}~\\
Notes & \cmd~\footnote{…}~ ; tableau : \cmd~threeparttable~, \cmd~\tnote{a}~, \w{tablenotes} ; crédit : \cmd~\source{…}~\\
\end{tbl}

\sect{7. Pièges courants}
\begin{tbl}{P{8.4cm}}
Compiler avec latexmk / LuaLaTeX, jamais pdfLaTeX.\\
\w{[?]} ou \w{??} dans le PDF : relancer (latexmk le fait seul).\\
Nombres : toujours le \textbf{point} dans le code, \cmd~\num{0.5}~ ; la sortie est en virgule.\\
Macros mathématiques : seulement en mode math (\cmd~$…$~).\\
Pas de ligne vide ni de \cmd~\code|…|~ dans les légendes.\\
\end{tbl}
\end{minipage}\hfill
%% ============================ COLONNE DROITE ============================
\begin{minipage}[t]{0.515\textwidth}
\cheatsize

\sect{8. Notations}
\begin{tbl}{l@{\enspace}l@{\qquad}l@{\enspace}l}
\cmd~\vect{v}~ & $\vect{v}$ & \cmd~\mat{A}~ & $\mat{A}$\\
\cmd~\uvect{n}~ & $\uvect{n}$ & \cmd~\tens{T}_{ij}~ & $\tens{T}_{ij}$\\
\cmd~\op{H}~ & $\op{H}$ & \cmd~\op{a}\adj~ & $\op{a}\adj$\\
\cmd~z\cc~ & $z\cc$ & \cmd~\mat{A}\transpose~ & $\mat{A}\transpose$\\
\cmd~\ii \ee \dd~ & $\ii\ \ee\ \dd$ & \cmd~\RR \CC \NN~ & $\RR\ \CC\ \NN$\\
\cmd~\ex \ey \ez~ & $\ex\ \ey\ \ez$ & \cmd~\er \etheta \ephi~ & $\er\ \etheta\ \ephi$\\
\end{tbl}

\sect{9. Dérivées et opérateurs différentiels}
\begin{tbl}{l@{\enspace}l@{\qquad}l@{\enspace}l}
\cmd~\dv{f}{x}~ & $\dv{f}{x}$ & \cmd~\dv[2]{f}{x}~ & $\dv[2]{f}{x}$\\
\cmd~\pdv{f}{x}~ & $\pdv{f}{x}$ & \cmd~\pdv[2]{f}{x}~ & $\pdv[2]{f}{x}$\\
\cmd~\pdv{f}{x}{y}~ & $\pdv{f}{x}{y}$ & \cmd~\fdv{S}{q}~ & $\fdv{S}{q}$\\
\cmd~\grad f~ & $\grad f$ & \cmd~\lapl f~ & $\lapl f$\\
\cmd~\diverg\vect{E}~ & $\diverg\vect{E}$ & \cmd~\curl\vect{B}~ & $\curl\vect{B}$\\
\end{tbl}

\sect{10. Mécanique quantique}
\begin{tbl}{l@{\enspace}l@{\qquad}l@{\enspace}l}
\cmd~\ket{\psi}~ & $\ket{\psi}$ & \cmd~\bra{\psi}~ & $\bra{\psi}$\\
\cmd~\braket{a}{b}~ & $\braket{a}{b}$ & \cmd~\ketbra{a}{b}~ & $\ketbra{a}{b}$\\
\cmd~\comm{A}{B}~ & $\comm{A}{B}$ & \cmd~\acomm{A}{B}~ & $\acomm{A}{B}$\\
\cmd~\poisson{f}{g}~ & $\poisson{f}{g}$ & \cmd~\pscal{a}{b}~ & $\pscal{a}{b}$\\
\cmd~\ev{\op{A}}~ & $\ev{\op{A}}$ & \cmd~\ev{\op{A}}{\psi}~ & $\ev{\op{A}}{\psi}$\\
\end{tbl}
\begin{tbl}{l@{\enspace}l}
\cmd~\mel{m}{\op{A}}{n}~ & $\mel{m}{\op{A}}{n}$\qquad \cmd~\ev*{…}{…}~ : taille automatique\\
\end{tbl}

\sect{11. Divers mathématiques}
\begin{tbl}{l@{\enspace}l@{\qquad}l@{\enspace}l}
\cmd~\abs{x}~ & $\abs{x}$ & \cmd~\abs*{\frac{a}{b}}~ & $\abs*{\frac{a}{b}}$\\
\cmd~\norm{\vect{v}}~ & $\norm{\vect{v}}$ & \cmd~\abs[\big]{x}~ & $\abs[\big]{x}$\\
\cmd~a\dotp b~ & $a\dotp b$ & \cmd~a\crossp b~ & $a\crossp b$\\
\cmd~\Tr \sgn \diag~ & $\Tr\ \sgn\ \diag$ & \cmd~\Reel z \Imag z~ & $\Reel z\ \Imag z$\\
\cmd~\bigO{x^3}~ & $\bigO{x^{3}}$ & \cmd~\hc~ & $\hc$\\
\end{tbl}
\par\vspace{0.1em}\cheatsize amsmath : \w{align gather cases} ; \cmd~\notag~ supprime un numéro.

\sect{12. Unités et incertitudes}
\begin{tbl}{l@{\enspace}l}
\cmd~\qty{9.81}{m.s^{-2}}~ & \qty{9.81}{m.s^{-2}}\\
\cmd~\qtyerr{9.75}{0.06}{m.s^{-2}}~ & \qtyerr{9.75}{0.06}{m.s^{-2}}\\
\cmd~\qtyrange{1}{5}{nm}~ & \qtyrange{1}{5}{nm}\\
\cmd~\num{1.2e-3}~ \ \cmd~\ang{30}~ & \num{1.2e-3}\quad\ang{30}\\
\cmd~\unit{\wavenumber}~ \ \cmd~\unit{eV}~ & \unit{\wavenumber}\quad\unit{eV}\\
\cmd~\incert{g}~ \ \cmd~\incertrel{g}~ & $\incert{g}$\quad $\incertrel{g}$\\
\cmd~\ecartnorm{a}{b}~ & $\ecartnorm{a}{b}$\\
\end{tbl}

\sect{13. Compiler (dossier du projet)}
\begin{tbl}{l@{\ }P{5.0cm}}
\cmd~compiler.bat~ & compile \w{main.tex}, vérifie, ouvre le PDF\\
\cmd~compiler.bat exemple-B-tp~ & compile un autre fichier\\
\cmd~compiler.bat nettoyer~ & supprime les fichiers auxiliaires\\
\cmd~latexmk fichier.tex~ & équivalent en ligne de commande\\
\end{tbl}
\end{minipage}

\end{document}
